% Veridian -- Workspace & Sync Architecture Design
% Compiled with: xelatex design.tex
% ----------------------------------------------------------------------------

\documentclass[12pt, a4paper]{article}

\usepackage{fontspec}
\usepackage{xeCJK}
\usepackage{geometry}
\usepackage{hyperref}
\usepackage{listings}
\usepackage{listingsutf8}
\usepackage{xcolor}
\usepackage{amssymb}
\usepackage{booktabs}
\usepackage{longtable}
\usepackage{array}
\usepackage{enumitem}
\usepackage{titlesec}
\usepackage{fancyhdr}
\usepackage{tcolorbox}
\usepackage{multirow}
\usepackage{tikz}
\usetikzlibrary{shapes.geometric, arrows.meta, positioning, fit, backgrounds, calc}
\tcbuselibrary{skins, breakable}

% ── Fonts ─────────────────────────────────────────────────────────────────────
\setmainfont{Times New Roman}
\setsansfont{Arial}
\setmonofont{Consolas}
\setCJKmainfont[BoldFont=SimHei, ItalicFont=KaiTi]{SimSun}
\setCJKsansfont{Microsoft YaHei}
\setCJKmonofont{FangSong}

% ── Page geometry ─────────────────────────────────────────────────────────────
\geometry{
  a4paper,
  top=25mm, bottom=25mm,
  left=28mm, right=28mm,
  headheight=24.5pt
}

% ── Colors ────────────────────────────────────────────────────────────────────
\definecolor{tsinghuapurple}{HTML}{660874}
\definecolor{codebg}{HTML}{F7F7F9}
\definecolor{codeframe}{HTML}{CCCCCC}
\definecolor{linkblue}{HTML}{0066CC}
\definecolor{sectionblue}{HTML}{003366}
\definecolor{ctrlplane}{HTML}{E8F4FD}
\definecolor{dataplane}{HTML}{FDF6E8}
\definecolor{clientlayer}{HTML}{F3E8FD}
\definecolor{warnbg}{HTML}{FFF8E1}
\definecolor{warnframe}{HTML}{F59E0B}
\definecolor{goodbg}{HTML}{F0FDF4}
\definecolor{goodframe}{HTML}{16A34A}
\definecolor{badbg}{HTML}{FEF2F2}
\definecolor{badframe}{HTML}{DC2626}

% ── Hyperref ──────────────────────────────────────────────────────────────────
\hypersetup{
  colorlinks=true,
  linkcolor=tsinghuapurple,
  urlcolor=linkblue,
  citecolor=tsinghuapurple,
  pdftitle={Veridian -- Workspace \& Sync Architecture Design},
  pdfauthor={RefAI Development Team},
}

% ── Section formatting ────────────────────────────────────────────────────────
\titleformat{\section}
  {\large\bfseries\color{sectionblue}}
  {\thesection}{1em}{}[{\color{tsinghuapurple}\titlerule[0.5pt]}]
\titleformat{\subsection}
  {\normalsize\bfseries\color{sectionblue}}
  {\thesubsection}{1em}{}
\titleformat{\subsubsection}
  {\normalsize\bfseries}
  {\thesubsubsection}{1em}{}

% ── tcolorbox styles ──────────────────────────────────────────────────────────
\tcbset{
  principlebox/.style={colback=tsinghuapurple!5, colframe=tsinghuapurple,
    fonttitle=\bfseries\small, left=6pt, right=6pt, top=4pt, bottom=4pt, breakable},
  fixbox/.style={colback=goodbg, colframe=goodframe,
    fonttitle=\bfseries\small, left=6pt, right=6pt, top=4pt, bottom=4pt, breakable},
  problembox/.style={colback=badbg, colframe=badframe,
    fonttitle=\bfseries\small, left=6pt, right=6pt, top=4pt, bottom=4pt, breakable},
  notebox/.style={colback=warnbg, colframe=warnframe,
    fonttitle=\bfseries\small, left=6pt, right=6pt, top=4pt, bottom=4pt, breakable},
}

% ── Code listings ─────────────────────────────────────────────────────────────
\lstset{
  basicstyle=\ttfamily\footnotesize, backgroundcolor=\color{codebg},
  frame=single, rulecolor=\color{codeframe}, breaklines=true,
  columns=fullflexible, keepspaces=true, showstringspaces=false, tabsize=2,
  xleftmargin=6pt, xrightmargin=6pt, aboveskip=4pt, belowskip=4pt,
  numbers=left, numberstyle=\tiny\color{gray}, numbersep=5pt,
  keywordstyle=\color{tsinghuapurple}\bfseries, commentstyle=\color{gray}\itshape,
  stringstyle=\color{linkblue},
}

% ── Header / Footer ───────────────────────────────────────────────────────────
\pagestyle{fancy}
\fancyhf{}
\fancyhead[L]{\small\color{gray} Veridian -- Workspace \& Sync Design}
\fancyhead[R]{\small\color{gray} 2026-07-08}
\fancyfoot[C]{\small\thepage}
\renewcommand{\headrulewidth}{0.4pt}
\renewcommand{\headrule}{\hbox to\headwidth{\color{tsinghuapurple}\leaders\hrule height \headrulewidth\hfill}}

\begin{document}

% ── Title page ────────────────────────────────────────────────────────────────
\begin{titlepage}
  \centering
  \vspace*{2.5cm}
  {\Huge\bfseries\color{tsinghuapurple} Veridian}\par
  \vspace{0.6em}
  {\Large\color{sectionblue} 工作空间与数据同步架构设计}\par
  \vspace{0.4em}
  {\large Workspace \& Data Synchronization Architecture Design}\par
  \vspace{2em}
  \hrule height 1.5pt \relax
  \vspace{2em}
  {\normalsize 控制面/数据面分离 \,·\, 可插拔同步后端 \,·\, 基于操作日志的合并 \,·\, 零强制云依赖}\par
  \vspace{3cm}
  {\small
  本方案在现有 Veridian 本地优先架构（\S 见 readme/refactor/redesign.tex）基础上，\\
  设计 Notion 式多人协作工作空间（邀请制、角色权限）与两种免费数据同步方式\\
  （GitHub 仓库 / 云盘文件夹），并逐条对照软件设计七大基本原则给出实现依据。
  }\par
  \vfill
  {\small\color{gray} Last updated: 2026-07-08 \quad 前置文档：readme/architecture/ · readme/refactor/}
\end{titlepage}

\tableofcontents
\newpage

% ═════════════════════════════════════════════════════════════════════════════
\section{背景与目标}
% ═════════════════════════════════════════════════════════════════════════════

\subsection{需求原文}

\begin{itemize}[leftmargin=2em, itemsep=2pt]
  \item 实现 Notion 式工作空间：可新建\textbf{邀请制多人协作}工作空间（邀请他人、分配权限），也可新建\textbf{私人工作空间}（个人多设备同步）。
  \item 数据同步方式一：在 GitHub 上创建仓库，用于团队共同管理和访问。
  \item 数据同步方式二：免费、无需承担费用、操作方便的备用同步方式。
  \item 权限模型：\textbf{自建轻量控制面服务}管理账号、成员、角色（而非完全委托给同步后端原生权限）。
\end{itemize}

\subsection{v1 范围声明}

\begin{center}
\begin{tabular}{p{6.5cm} p{6.5cm}}
  \toprule
  \textbf{本次设计包含} & \textbf{明确不含（见 \S 12 Future Work）} \\
  \midrule
  邀请制工作空间 + 角色权限     & 真正对外公开可发现的工作空间 \\
  GitHub 仓库同步后端           & 实时协作光标 / 同时编辑同一字段的即时广播 \\
  云盘文件夹同步后端             & 端到端加密（本设计预留密钥管理钩子，v1 默认不开启） \\
  基于 oplog 的异步合并 + 冲突提示 & WebDAV / Syncthing 后端（架构可扩展，暂不实现） \\
  \bottomrule
\end{tabular}
\end{center}

\begin{tcolorbox}[notebox, title={与现有重构路线图的关系}]
  本设计是《readme/refactor/redesign.tex》中 \textbf{P4（同步就绪）} 阶段预埋的
  版本号、墓碑表（tombstones）、操作日志（oplog）三项机制的\textbf{直接消费方}。
  这些机制在 P2 重构时就已经写入 \texttt{db/migrations} 与 \texttt{Notifier} 事件总线，
  本次不再重新设计数据模型，只是把"谁来读/写 oplog"这件事从"本地进程自己"
  扩展为"多台设备 + 多个协作者"。
\end{tcolorbox}

% ═════════════════════════════════════════════════════════════════════════════
\section{总体架构：控制面与数据面分离}
% ═════════════════════════════════════════════════════════════════════════════

这是本设计的核心决策。把系统拆成两个职责完全独立、可分别替换的平面：

\begin{center}
\begin{tabular}{p{3cm} p{4.6cm} p{4.6cm}}
  \toprule
  & \textbf{控制面 (Control Plane)} & \textbf{数据面 (Data Plane)} \\
  \midrule
  职责       & 身份、工作空间成员、角色、邀请令牌 & 实际文献数据（oplog + 附件二进制） \\
  存储位置   & 自建后端（推荐 Supabase 免费层） & GitHub 仓库 \textbf{或} 用户自己的云盘文件夹 \\
  数据量     & 极小（几十 KB 级的成员/角色表）  & 大（PDF、Markdown、图片，GB 级） \\
  是否收费   & 免费层足够（见 \S 11 成本表）    & 完全免费——复用用户已有的 GitHub/云盘账号 \\
  单点故障风险 & 有（但数据量小，易备份/迁移）    & 无——数据本体始终在用户自己的仓库/网盘里 \\
  \bottomrule
\end{tabular}
\end{center}

\begin{tcolorbox}[principlebox, title={为什么要分离}]
  控制面只回答"你是谁、你在哪些工作空间、你有什么权限"这类\textbf{小数据、强一致性}问题，
  适合放在一个真正的关系型数据库里做行级安全策略（RLS）。
  数据面只回答"字节怎么从 A 设备搬到 B 设备"这类\textbf{大数据、最终一致性}问题，
  适合交给用户已有的免费存储（GitHub/Google Drive/OneDrive）。
  两者用不同的一致性模型、不同的容量规划、不同的失败模式——
  混在一起会让 GitHub 仓库既要当"数据库"又要当"文件仓库"，两头都做不好。
\end{tcolorbox}

\subsection{分层架构图}

\begin{center}
\begin{tikzpicture}[
  layer/.style={rectangle, rounded corners=5pt, draw=sectionblue, thick,
                minimum width=13cm, minimum height=1.3cm, font=\small, align=center},
  sub/.style={rectangle, rounded corners=3pt, draw=gray!60,
              minimum height=0.7cm, font=\footnotesize, align=center, fill=white},
  arr/.style={-Stealth, very thick, color=tsinghuapurple},
]
  % Layer: UI
  \node[layer, fill=clientlayer] (ui) {};
  \node[anchor=west, font=\small\bfseries\color{sectionblue}] at ($(ui.west)+(0.2,0.35)$) {客户端 UI 层（渲染进程）};
  \node[sub, anchor=east] at ($(ui.east)+(-0.2,-0.15)$) (uisub3) {ConflictBanner};
  \node[sub, anchor=east] at ($(uisub3.west)+(-0.15,0)$) (uisub2) {MembersPanel / InviteDialog};
  \node[sub, anchor=east] at ($(uisub2.west)+(-0.15,0)$) (uisub1) {WorkspaceSwitcher};

  % Layer: client services
  \node[layer, fill=clientlayer!60, below=0.4cm of ui] (svc) {};
  \node[anchor=west, font=\small\bfseries\color{sectionblue}] at ($(svc.west)+(0.2,0.35)$) {客户端服务层（主进程，Electron）};
  \node[sub, anchor=east] at ($(svc.east)+(-0.2,-0.15)$) (svcsub3) {SyncEngine};
  \node[sub, anchor=east] at ($(svcsub3.west)+(-0.15,0)$) (svcsub2) {PermissionSource};
  \node[sub, anchor=east] at ($(svcsub2.west)+(-0.15,0)$) (svcsub1) {WorkspaceService};

  % Layer: control plane
  \node[layer, fill=ctrlplane, below=0.55cm of svc] (ctrl) {};
  \node[anchor=west, font=\small\bfseries\color{sectionblue}] at ($(ctrl.west)+(0.2,0.35)$) {控制面（自建后端 -- Supabase 免费层）};
  \node[sub, anchor=east] at ($(ctrl.east)+(-0.2,-0.15)$) (ctrlsub3) {Realtime（成员/角色变更推送）};
  \node[sub, anchor=east] at ($(ctrlsub3.west)+(-0.15,0)$) (ctrlsub2) {Postgres + RLS（workspaces/members/invites）};
  \node[sub, anchor=east] at ($(ctrlsub2.west)+(-0.15,0)$) (ctrlsub1) {Auth（邮箱/OAuth 登录）};

  % Layer: data plane
  \node[layer, fill=dataplane, below=0.55cm of ctrl] (data) {};
  \node[anchor=west, font=\small\bfseries\color{sectionblue}] at ($(data.west)+(0.2,0.35)$) {数据面（可插拔 SyncBackend）};
  \node[sub, anchor=east] at ($(data.east)+(-0.2,-0.15)$) (datasub2) {CloudFolderSyncBackend（Drive/OneDrive）};
  \node[sub, anchor=east] at ($(datasub2.west)+(-0.15,0)$) (datasub1) {GitSyncBackend（GitHub 仓库）};

  \draw[arr] (ui.south) -- (svc.north);
  \draw[arr] (svc.south) -- node[right, font=\tiny]{鉴权/角色查询} (ctrl.north);
  \draw[arr] (svc.south) -- node[left, font=\tiny]{oplog + 附件推拉} ($(data.north)+(-3.5,0)$);
\end{tikzpicture}
\end{center}

% ═════════════════════════════════════════════════════════════════════════════
\section{控制面设计：工作空间、成员、角色}
% ═════════════════════════════════════════════════════════════════════════════

\subsection{为什么选 Supabase}

\begin{center}
\begin{tabular}{p{3cm} p{9.5cm}}
  \toprule
  \textbf{维度} & \textbf{说明} \\
  \midrule
  数据模型 & Postgres 关系型数据库，工作空间/成员/角色天然是关系型数据，
             比 Firebase Firestore 的文档模型更贴合，也更接近 Veridian 本地
             SQLite 的建模习惯，schema 可以直接照抄改写 \\
  权限实现 & 内置\textbf{行级安全策略（Row Level Security）}，可以在数据库层
             再加一道"用户只能读自己所在工作空间的行"的防线，
             不完全依赖应用层代码 \\
  免费层   & 500 MB 数据库 + 5 万月活用户 + 内置 Auth + Realtime 推送，
             对于"成员表 + 角色表"这种量级的数据完全够用，个人/小团队
             永远用不满 \\
  账号体系 & 内置邮箱密码 / Magic Link / GitHub OAuth 登录，
             GitHub OAuth 登录还能顺便一次性拿到 GitHub 授权，
             给后面接 GitSyncBackend 省一步 \\
  \bottomrule
\end{tabular}
\end{center}

\subsection{控制面数据模型}

\begin{longtable}{p{3cm} p{9.5cm}}
  \toprule
  \textbf{表} & \textbf{字段说明} \\
  \midrule
  \texttt{workspaces} & id, name, kind (\texttt{'private'} | \texttt{'shared'}),
    owner\_id, sync\_backend\_type, sync\_backend\_config（JSON，仓库地址/云盘路径等
    非敏感连接信息；令牌类敏感信息不落库，见 \S 5.3）, created\_at \\
  \texttt{workspace\_members} & workspace\_id, user\_id, role
    (\texttt{owner}|\texttt{admin}|\texttt{editor}|\texttt{viewer}), joined\_at \\
  \texttt{invites} & id, workspace\_id, email, role, token（一次性）,
    expires\_at, status (\texttt{pending}|\texttt{accepted}|\texttt{revoked}) \\
  \texttt{devices} & id, user\_id, device\_name, last\_synced\_at
    （用于"私人工作空间"的多设备同步场景，展示"我的设备列表"） \\
  \bottomrule
\end{longtable}

\subsection{角色权限矩阵}

\begin{center}
\begin{tabular}{p{2.6cm} p{2.3cm} p{2.3cm} p{2.3cm} p{2.6cm}}
  \toprule
  \textbf{操作} & \textbf{Owner} & \textbf{Admin} & \textbf{Editor} & \textbf{Viewer} \\
  \midrule
  查看文献/附件         & \checkmark & \checkmark & \checkmark & \checkmark \\
  新增/编辑/删除文献     & \checkmark & \checkmark & \checkmark & $\times$ \\
  管理集合（文件夹）     & \checkmark & \checkmark & \checkmark & $\times$ \\
  邀请 Viewer/Editor    & \checkmark & \checkmark & $\times$   & $\times$ \\
  邀请 Admin，移除成员    & \checkmark & \checkmark & $\times$   & $\times$ \\
  修改同步后端配置       & \checkmark & $\times$   & $\times$   & $\times$ \\
  删除工作空间           & \checkmark & $\times$   & $\times$   & $\times$ \\
  \bottomrule
\end{tabular}
\end{center}

角色定义是一张静态的\textbf{角色 $\to$ 权限集合}映射表（见 \S 7 ISP/OCP 小节），
新增角色（如未来的"Commenter"）只需扩表，不需要改任何判断权限的调用点。

\subsection{邀请流程时序}

\begin{center}
\begin{tikzpicture}[
  actor/.style={font=\footnotesize\bfseries},
  msg/.style={-Stealth, thick},
  lifeline/.style={dashed, gray},
]
  \foreach \x/\name in {0/Owner 客户端, 4/Supabase 控制面, 8/被邀请人客户端} {
    \node[actor] at (\x, 0) {\name};
    \draw[lifeline] (\x,-0.3) -- (\x,-7);
  }
  \draw[msg] (0,-1) -- node[above, font=\tiny]{createInvite(email, role)} (4,-1);
  \draw[msg] (4,-1.8) -- node[above, font=\tiny]{生成一次性 token，发送邀请邮件} (4,-1.8);
  \draw[msg, dashed] (4,-2.6) -- node[above, font=\tiny]{（邮件外部投递）} (8,-2.6);
  \draw[msg] (8,-3.4) -- node[above, font=\tiny]{acceptInvite(token) + 登录/注册} (4,-3.4);
  \draw[msg] (4,-4.2) -- node[above, font=\tiny]{校验 token 未过期 → 写入 workspace\_members} (4,-4.2);
  \draw[msg] (4,-5) -- node[above, font=\tiny]{Realtime 推送成员变更} (0,-5);
  \draw[msg] (4,-5.8) -- node[above, font=\tiny]{返回 sync\_backend\_config（仓库地址/云盘路径）} (8,-5.8);
  \draw[msg] (8,-6.6) -- node[above, font=\tiny]{SyncEngine 首次 pull 全量 oplog} (8,-6.6);
\end{tikzpicture}
\end{center}

被邀请人客户端拿到的是\textbf{连接信息}（仓库地址、云盘路径），
而不是数据本体——数据本体在首次 pull 时才从数据面拉取，
控制面全程不经手任何文献数据。

% ═════════════════════════════════════════════════════════════════════════════
\section{数据面设计：可插拔同步后端}
% ═════════════════════════════════════════════════════════════════════════════

\subsection{SyncBackend 抽象}

\begin{lstlisting}[language={}]
// shared/sync/SyncBackend.ts -- 纯传输接口，不含任何权限概念
// （权限已经在 \S3 的控制面解决，后端只负责字节搬运）
export interface SyncBackend {
  testConnection(): Promise<boolean>

  // 拉取自 cursor（GitSyncBackend 用 commit SHA，
  // CloudFolderSyncBackend 用 manifest 版本号）之后的新变更包
  pull(cursor: string | null): Promise<{ envelopes: SyncEnvelope[]; nextCursor: string }>

  // 推送本地产生的变更包（内含 oplog 条目 + 引用的附件内容哈希）
  push(envelope: SyncEnvelope): Promise<{ ok: boolean; conflict?: boolean }>

  // 按内容哈希上传/下载附件二进制（去重：哈希已存在则跳过上传）
  uploadBlob(hash: string, data: Uint8Array): Promise<void>
  downloadBlob(hash: string): Promise<Uint8Array>
}

export interface SyncEnvelope {
  deviceId: string
  workspaceId: string
  seqRange: [number, number]     // 对应本地 oplog.seq 区间
  entries: OplogEntry[]          // 复用现有 oplog 表结构
  ts: number
}
\end{lstlisting}

\subsection{GitSyncBackend（GitHub 仓库）}

\begin{itemize}[leftmargin=2em, itemsep=2pt]
  \item 使用 \textbf{isomorphic-git}（纯 JS 实现）而非依赖系统 git 可执行文件——
        Electron 应用免于要求用户预装 Git，跨平台打包更简单。
  \item 仓库布局（内容寻址，天然无 git 合并冲突）：
\end{itemize}

\begin{lstlisting}[language={}, numbers=none]
<repo>/
├── oplog/<device-id>/<seq-start>-<seq-end>.jsonl   # 追加写入
├── blobs/<sha256前2位>/<sha256>                     # 附件二进制，按哈希去重
└── manifest.json                                    # 各设备已知的最新 seq（仅供快速跳过）
\end{lstlisting}

\begin{itemize}[leftmargin=2em, itemsep=2pt]
  \item PUSH = 写入新的 oplog 分片 + 未见过哈希的 blob 文件 $\to$ commit $\to$ push。
  \item PULL = fetch + fast-forward merge；由于文件按设备目录、内容哈希隔离，
        \textbf{结构上不会产生 git 合并冲突}（这是"CRDT via git"的常见模式）。
  \item 认证：GitHub Personal Access Token（细粒度 PAT，仅 \texttt{contents:write}
        权限），通过 Electron \texttt{safeStorage} 加密存储——复用现有 MinerU API
        Token 的加密落地模式（见 readme/refactor/redesign.tex \S3）。
  \item 大文件：单文件 $>$ 50 MB 时提示启用 Git LFS（GitHub 免费层：1 GB 存储 +
        1 GB/月带宽），避免仓库体积膨胀影响 clone 速度。
\end{itemize}

\subsection{CloudFolderSyncBackend（云盘文件夹）}

\begin{itemize}[leftmargin=2em, itemsep=2pt]
  \item \textbf{核心思路}：Veridian 不直接对接任何云盘 API，只是把上面同样的
        \texttt{oplog/ blobs/ manifest.json} 布局，写在用户指定的一个\textbf{本地目录}里——
        这个目录恰好是 Google Drive / OneDrive / Dropbox 桌面客户端正在同步的目录。
        真正的字节搬运完全交给用户已经装好、免费的云盘客户端，Veridian 只管本地读写。
  \item 用 \texttt{chokidar} 监听该目录，检测到\textbf{其他设备}写入的新 oplog 分片后触发 pull-replay。
  \item \textbf{零 API Key、零 OAuth、零第三方服务账号}——这是"免费、不承担费用、
        方便操作"这条需求里契合度最高的方案，因为几乎所有用户机器上都已经装了
        至少一种云盘客户端。
  \item 已知限制：SQLite 数据库文件本身\textbf{不可以}直接丢进云盘同步目录
        （云盘客户端可能在数据库写入过程中读到半写状态导致损坏）——这正是
        为什么要通过 oplog 分片文件（一次性追加写完才算存在）而不是同步
        \texttt{.db} 文件本身。
\end{itemize}

\begin{tcolorbox}[fixbox, title={对比：为什么两个后端接口能长得一模一样}]
  GitSyncBackend 和 CloudFolderSyncBackend 底层设施完全不同（Git 协议 vs 文件系统 +
  第三方客户端），但对上层 \texttt{SyncEngine} 暴露的接口完全一致——
  这正是 \S 7 里\textbf{里氏替换原则（LSP）}的直接体现：换后端不需要改一行
  \texttt{SyncEngine} 代码。
\end{tcolorbox}

\subsection{冲突解决策略（v1：最后写入优先 + 手动提示）}

\begin{enumerate}[leftmargin=2em, itemsep=2pt]
  \item Pull 到的 oplog 条目通过\textbf{与本地写入完全相同的 Service 层函数}回放
        （\texttt{ItemService.updateItem} 等），复用 P2 重构里"单一写入口必发事件"
        的设计——远程变更到达后 UI 自动刷新，无需为同步单独写一套 UI 逻辑。
  \item 字段级最后写入优先：比较 \texttt{items.version} + 操作时间戳，
        新版本覆盖旧版本（与 Zotero 自身的合并策略一致）。
  \item 真正的并发冲突（两台设备在离线期间都改了同一文献的同一字段）
        会被识别（版本号分叉）并放入 ConflictBanner，用户手动选择保留哪一份，
        而不是静默丢弃——这是 v1 里唯一需要人工介入的场景。
\end{enumerate}

% ═════════════════════════════════════════════════════════════════════════════
\section{权限与安全}
% ═════════════════════════════════════════════════════════════════════════════

\subsection{PermissionSource 抽象}

\begin{lstlisting}[language={}]
// main/services/PermissionSource.ts
export interface PermissionSource {
  getRole(workspaceId: string, userId: string): Promise<Role | null>
  subscribeRoleChanges(workspaceId: string, cb: (change: RoleChange) => void): () => void
}

// 目前唯一实现：SupabasePermissionSource（查 workspace_members 表 +
// 订阅 Supabase Realtime）。未来若要支持"自托管控制面"，
// 只需新增一个实现，WorkspaceService 不用改。
\end{lstlisting}

\subsection{令牌与密钥的存储原则}

\begin{center}
\begin{tabular}{p{3.5cm} p{4.6cm} p{4.6cm}}
  \toprule
  \textbf{凭据} & \textbf{存储位置} & \textbf{加密方式} \\
  \midrule
  Supabase 会话 token   & 本地（Electron userData） & \texttt{safeStorage} 加密 \\
  GitHub PAT           & 本地（每设备各自持有）    & \texttt{safeStorage} 加密，\textbf{不}同步到控制面 \\
  云盘文件夹路径         & 本地设置                  & 明文（非敏感） \\
  \bottomrule
\end{tabular}
\end{center}

\begin{tcolorbox}[notebox, title={端到端加密（Future Work，非 v1 必须）}]
  v1 的 oplog/blob 内容在数据面（GitHub/云盘）中以明文存储，
  安全边界依赖"用户信任自己选择的仓库/网盘"。
  若未来需要更强隐私保证，可在 \texttt{workspaces} 表加一列
  \texttt{encrypted\_workspace\_key}（用每个成员的公钥做信封加密），
  \texttt{SyncEnvelope} 序列化前先用该对称密钥加密——
  控制面和数据面提供方都只能看到密文。本设计已把该密钥字段的位置预留好，
  但不在 v1 实现范围内。
\end{tcolorbox}

% ═════════════════════════════════════════════════════════════════════════════
\section{客户端架构改动清单}
% ═════════════════════════════════════════════════════════════════════════════

\begin{longtable}{p{4.4cm} p{8.1cm}}
  \toprule
  \textbf{新增模块} & \textbf{职责} \\
  \midrule
  \texttt{main/services/WorkspaceService.ts} & 工作空间 CRUD、成员邀请/移除、当前工作空间切换；
    唯一的写入口，调用后发 \texttt{workspace.changed} 域事件 \\
  \texttt{main/services/PermissionSource.ts} & 角色查询抽象 + Supabase 实现（\S5.1） \\
  \texttt{main/sync/SyncEngine.ts} & 编排 pull/push 时机、调用 \texttt{SyncBackend}、
    驱动 oplog 回放、检测冲突 \\
  \texttt{main/sync/backends/GitSyncBackend.ts} & \S4.2 实现 \\
  \texttt{main/sync/backends/CloudFolderSyncBackend.ts} & \S4.3 实现 \\
  \texttt{main/sync/OplogReplayer.ts} & 把远程 oplog 条目映射回 Service 层调用 \\
  \texttt{shared/sync/SyncBackend.ts} & 传输接口定义（\S4.1），主/渲染进程共享类型 \\
  \texttt{renderer/.../WorkspaceSwitcher.tsx} & 顶部工具栏的工作空间切换下拉 \\
  \texttt{renderer/.../InviteDialog.tsx} & 输入邮箱 + 选择角色 $\to$ 调用 WorkspaceService \\
  \texttt{renderer/.../MembersPanel.tsx} & 成员列表 + 角色调整 + 移除（按权限矩阵置灰） \\
  \texttt{renderer/.../ConflictBanner.tsx} & 展示待手动解决的字段级冲突 \\
  \bottomrule
\end{longtable}

新增 IPC 通道遵循现有\texttt{shared/ipc-contract.ts}契约化模式（zod schema + 网关校验），
不额外发明新的 IPC 约定。

% ═════════════════════════════════════════════════════════════════════════════
\section{七大设计原则逐一对应}
% ═════════════════════════════════════════════════════════════════════════════

\begin{longtable}{p{2.6cm} p{9.9cm}}
  \toprule
  \textbf{原则} & \textbf{在本设计中的体现} \\
  \midrule
  单一职责原则 \newline (SRP) &
    \texttt{SyncEngine}（何时同步、冲突判定）与 \texttt{SyncBackend}（字节怎么搬）
    职责完全分离；\texttt{WorkspaceService}（成员/角色管理）与
    \texttt{PermissionSource}（角色查询）分离；\texttt{OplogReplayer}
    （把远程条目映射回本地调用）单独成一个模块，不与传输逻辑混在一起。
    每个模块只有一个改变它的理由。 \\
  \midrule
  开闭原则 \newline (OCP) &
    新增同步后端（WebDAV、Syncthing）只需新写一个实现
    \texttt{SyncBackend} 接口的类，\texttt{SyncEngine}/UI 不需要改一行；
    新增角色只需在角色$\to$权限集合的静态映射表里加一行（\S3.3），
    所有权限判断代码天然生效，不需要改判断逻辑本身。 \\
  \midrule
  里氏替换原则 \newline (LSP) &
    \texttt{GitSyncBackend} 与 \texttt{CloudFolderSyncBackend} 对
    \texttt{SyncEngine} 而言完全可互换（\S4.3 fixbox）：两者都保证
    push 要么完整成功要么完整失败（不会半推送），pull 不会丢失条目——
    这两条前置/后置条件对所有实现都一致成立。 \\
  \midrule
  接口隔离原则 \newline (ISP) &
    \texttt{SyncBackend} 只包含传输方法（pull/push/blob 存取），
    不包含任何权限判断方法——权限完全隔离进独立的 \texttt{PermissionSource}
    接口，两个接口的消费者（\texttt{SyncEngine} vs \texttt{WorkspaceService}）
    互不依赖对方用不到的方法。 \\
  \midrule
  依赖倒置原则 \newline (DIP) &
    \texttt{SyncEngine} 和 \texttt{WorkspaceService} 依赖的是
    \texttt{SyncBackend}/\texttt{PermissionSource} 这两个抽象接口，
    具体用哪个实现（Git/CloudFolder，Supabase/未来的自托管方案）
    由工作空间配置在运行时注入，高层模块不 import 任何具体后端类。 \\
  \midrule
  迪米特法则 \newline (LoD，最少知识原则) &
    UI 组件只通过 \texttt{WorkspaceService} 的 IPC 契约通信，
    从不直接触碰 \texttt{SyncBackend} 内部（GitHub API 细节、云盘路径拼接）
    或 oplog 表结构——延续 P2 重构"单一写入口、薄 IPC 处理器"的既有纪律。 \\
  \midrule
  合成复用原则 \newline (CARP，组合优于继承) &
    \texttt{SyncEngine} 是通过\textbf{组合}一个可替换的 \texttt{SyncBackend} +
    \texttt{OplogReplayer} + 冲突判定逻辑构成的，而不是为每种后端派生一个
    \texttt{GitSyncEngine extends BaseSyncEngine} 子类——避免脆弱基类问题，
    每个组件可独立单元测试（呼应 redesign.tex \S9 的测试金字塔目标）。 \\
  \bottomrule
\end{longtable}

% ═════════════════════════════════════════════════════════════════════════════
\section{与现有重构路线图的衔接}
% ═════════════════════════════════════════════════════════════════════════════

\begin{center}
\begin{tabular}{p{2.2cm} p{4.4cm} p{6.1cm}}
  \toprule
  \textbf{阶段} & \textbf{redesign.tex 中的产物} & \textbf{本设计如何复用} \\
  \midrule
  P2 & Notifier 事件总线 + Service 层 & OplogReplayer 复用同一批 Service 函数回放远程变更，
       远程变更自动触发 UI 刷新，零新增前端状态管理代码 \\
  P3 & DB Worker 线程 + UnitOfWork 事务 & SyncEngine 的 pull 回放批量操作复用同一套事务封装，
       避免半回放状态 \\
  P4 & 版本号 + 墓碑表 + oplog & 本设计的数据面协议直接就是"如何搬运这张 oplog 表"，
       零数据模型变更 \\
  P5 & 托管存储布局 \texttt{storage/<item-key>/} & 附件内容寻址哈希直接对该目录下的文件计算，
       天然对齐 \\
  \bottomrule
\end{tabular}
\end{center}

% ═════════════════════════════════════════════════════════════════════════════
\section{分阶段实施路线图}
% ═════════════════════════════════════════════════════════════════════════════

\begin{center}
\begin{tabular}{p{1.4cm} p{5.6cm} p{2.4cm} p{2.6cm}}
  \toprule
  \textbf{阶段} & \textbf{内容} & \textbf{工作量} & \textbf{风险} \\
  \midrule
  W1 & Supabase 项目搭建：schema + RLS + Auth 接入；
       客户端登录/注册 UI & 3--4 天 & 低 \\
  W2 & WorkspaceService + PermissionSource + 工作空间切换 UI（暂不接同步，
       先验证控制面闭环） & 4--5 天 & 低 \\
  W3 & SyncBackend 抽象 + CloudFolderSyncBackend（先做这个，
       零外部依赖，验证 oplog 传输协议） & 4--5 天 & 中 \\
  W4 & GitSyncBackend（isomorphic-git 接入 + PAT 加密存储 + LFS 提示） & 5--6 天 & 中 \\
  W5 & 冲突检测 + ConflictBanner UI + 邀请流程端到端联调 & 4--5 天 & 中 \\
  测试 & 每阶段同步编写：SyncBackend 用假后端跑单测，
        OplogReplayer 用内存库跑集成测试 & 贯穿 & -- \\
  \bottomrule
\end{tabular}
\end{center}

约 4 周完成 v1。W1--W2（控制面）与 W3--W4（数据面）在依赖上是并行的，
如有两位开发者可并行推进以压缩到约 2.5 周。

% ═════════════════════════════════════════════════════════════════════════════
\section{成本与限制}
% ═════════════════════════════════════════════════════════════════════════════

\begin{longtable}{p{3.2cm} p{4cm} p{5.1cm}}
  \toprule
  \textbf{服务} & \textbf{免费额度} & \textbf{对个人/小团队库的估算} \\
  \midrule
  Supabase & 500 MB 数据库，5 万 MAU，
    2 GB 文件存储（本设计不用它存文献数据） & 一个工作空间的成员/角色/邀请记录
    通常 $<$ 1 MB；千人协作规模才会接近上限 \\
  GitHub（私有仓库） & 免费账号：无限私有仓库、无限协作者；
    LFS 免费层 1 GB 存储 + 1 GB/月带宽 & 纯文本（oplog + markdown）几乎不占额度；
    PDF 附件量大时需要开 LFS 或改用云盘后端 \\
  Google Drive / OneDrive & 15 GB / 5 GB 免费额度 & 覆盖大多数个人文献库；
    团队场景建议每个协作者用自己的网盘额度存自己上传的原文件 \\
  \bottomrule
\end{longtable}

\begin{tcolorbox}[notebox, title={何时需要考虑付费}]
  单个工作空间协作者超过几百人，或单个云盘账号存储的文献库超过其免费额度时，
  才需要考虑：GitHub LFS 付费带宽包、云盘更高存储层级，或自建 S3 兼容存储
  作为数据面第三种后端（架构上只需新增一个 \texttt{SyncBackend} 实现）。
  对个人和小型科研团队场景，v1 设计\textbf{完全免费}。
\end{tcolorbox}

% ═════════════════════════════════════════════════════════════════════════════
\section{待扩展功能（Future Work）}
% ═════════════════════════════════════════════════════════════════════════════

\begin{itemize}[leftmargin=2em, itemsep=2pt]
  \item 端到端加密（\S5.2 已预留密钥字段）
  \item 真正对外公开可发现的工作空间（类似开源项目 + PR 审核合并模式）
  \item 实时协作（同一文献同时编辑时的即时光标广播，需要 WebSocket/CRDT，
        当前 oplog 异步合并模型不支持）
  \item WebDAV / Syncthing 作为第三、第四种 \texttt{SyncBackend} 实现
  \item 自托管控制面（\texttt{PermissionSource} 第二实现，替代 Supabase 依赖）
  \item 移动端 / Web 端客户端接入同一套控制面 + 数据面协议
\end{itemize}

\end{document}
